% =====================================================================
%  2. Contexto del problema
% =====================================================================
\section{Contexto del problema}
\label{sec:contexto}

El problema que aborda esta investigación se sitúa en la intersección de tres
hechos documentados: la carga de trabajo docente y su composición real, la
definición de calidad de un recurso educativo y los límites de lo que hoy se
automatiza, y el estado de la investigación sobre IA en educación superior.

\subsection{La carga docente y sus componentes invisibles}
\label{sec:contexto-carga}

La preparación y el mantenimiento de material de curso forman parte de un conjunto
de actividades que los modelos institucionales de asignación de carga tienden a
subestimar. \textcite{kenny2023life} validaron una herramienta holística de
estimación de carga académica (AWET) con 39 académicos de 21 universidades
australianas, contrastando la carga estimada con la asignada nominalmente. La carga
promedio ponderada resultó del \textbf{132\,\% de una carga completa} (unas
2252~horas anuales); 24 de los 39 participantes superaban el 110\,\% y 11 superaban
el 150\,\%. Solo 3 tenían un perfil efectivamente equilibrado entre docencia,
investigación y servicio, frente a 18 según su asignación nominal. Las actividades
de coordinación, liderazgo, administración, supervisión y adaptación de la
enseñanza en línea no estaban contempladas en los modelos institucionales.

\begin{advertencia}
\textcite{kenny2023life} no miden la carga específica de curar o mantener
documentación de cursos, y su \enquote{validación} es cualitativa: no reportan
indicadores psicométricos. La fuente sirve para argumentar que existe trabajo
docente de apoyo sistemáticamente no contabilizado, \emph{no} para cuantificar el
tiempo que un docente dedica a la curación.
\end{advertencia}

\subsection{La expectativa sobre la IA generativa y lo que la evidencia muestra}
\label{sec:contexto-genai}

La hipótesis implícita de que automatizar equivale a reducir trabajo no se sostiene
con la evidencia disponible. La revisión sistemática de \textcite{su2026efficiency}
--- 30 estudios empíricos, cribado doble con $\kappa = 0{,}845$ en título y resumen
y $\kappa = 0{,}851$ en texto completo, calidad evaluada con MMAT~2018 ---
documenta una redistribución: la IA generativa ahorra tiempo en preparación de
materiales, borradores de retroalimentación y textos administrativos, pero añade
revisión de la calidad de las salidas, verificación de hechos, juicio ético y
monitoreo de integridad académica.

Un hallazgo de esa revisión condiciona directamente el diseño de evaluación de este
proyecto: las conclusiones de los estudios primarios dependen de \emph{cómo} midan
la carga. Los que la operacionalizan como eficiencia (tiempo de tarea) tienden a
encontrar reducción; los que la operacionalizan como intensificación (adaptación,
supervisión, responsabilidad) tienden a encontrar aumento. Los autores advierten
que \textcquote[p.~10]{su2026efficiency}{Future research that relies only
on time-efficiency indicators may underestimate the invisible work intensification
brought about by GenAI}. De ahí que la Parte~V contemple tiempo \emph{y} carga
percibida como medidas complementarias, no intercambiables.

\subsection{Qué significa que un recurso educativo tenga calidad}
\label{sec:contexto-calidad}

La calidad de un recurso educativo es multidimensional, pero la literatura converge
en que la exactitud del contenido es la dimensión dominante. El
\emph{Learning Object Review Instrument} (LORI~1.5), cuyos fundamentos expone
\textcite{leacock2007framework}, define nueve dimensiones --- calidad del contenido,
alineación con objetivos, retroalimentación y adaptación, motivación, diseño de la
presentación, usabilidad, accesibilidad, reusabilidad y cumplimiento de estándares ---
y sitúa la primera por encima de las demás:
\textcquote[p.~45]{leacock2007framework}{A learning resource is of little
or no use if it is well designed in all other respects but its content is inaccurate
or misleading}. Los autores señalan además que los estudiantes desarrollan
comprensiones incompletas que reproducen los errores de los materiales, y que la
presencia de información específica de una instancia --- fechas de entrega, nombres
del instructor --- degrada la reusabilidad de un recurso; ambos puntos anticipan dos
de las clases de defecto que este proyecto busca detectar: la contradicción y la
desactualización.

El estudio de \textcite{xie2018systematic} sobre 1200 reseñas de recursos digitales
elaboradas por docentes aporta un dato metodológicamente relevante: los revisores
comentan con más frecuencia las virtudes que los defectos de un recurso. La revisión
humana, por sí sola, tiende a subdetectar problemas. En sentido complementario,
\textcite{tabares2017modelo}, al aplicar un modelo de evaluación por capas a 46
objetos de aprendizaje de la Federación de Repositorios de Objetos de Aprendizaje de
Colombia, documentan el caso inverso: un objeto con buenos indicadores automáticos
--- basados en metadatos --- obtuvo valoraciones muy bajas de expertos y usuarios.
Ninguna de las dos vías basta por separado.

\begin{estado}
\small
\textbf{Lectura conjunta de las tres fuentes.} La evaluación automática detecta lo
que es formalizable (completitud de metadatos, enlaces rotos, similitud entre
textos); la revisión humana juzga la sustancia pero subdetecta defectos y no escala.
La combinación de ambas es la posición que esta investigación adopta, y coincide con
la recomendación de \textcite{clements2015open} tras revisar 82 artículos sobre
repositorios educativos: \textcquote[p.~6]{clements2015open}{Expert review
might be not the most economical approach, but it seems to be needed in order to
evaluate the substance of the resources in the repository}.
\end{estado}

\subsection{Qué automatiza hoy el aseguramiento de calidad}
\label{sec:contexto-automatizacion}

La revisión de \textcite{clements2015open} identifica quince instrumentos
específicos de aseguramiento de calidad en repositorios educativos: revisión por
pares, calificaciones, herramientas tipo LORI, recomendadores, comentarios,
etiquetado social, marcado de problemas (enlaces rotos, contenido inapropiado),
especificación de autoría y licencias, redes de confianza y pruebas automáticas de
metadatos. El alcance de lo automatizado en ese panorama se limita a
\textbf{metadatos, enlaces y recomendaciones}: ningún instrumento verifica el
contenido en sí. Los autores observan además que los instrumentos colaborativos
dependen de una comunidad activa y que, en educación, la motivación para calificar y
comentar es menor que en comercio electrónico --- una restricción que se agrava en
el corpus de un curso individual, donde la comunidad es, en el límite, una persona.

\subsection{El estado de la investigación sobre IA en educación superior}
\label{sec:contexto-aied}

Dos revisiones separadas por cinco años coinciden en dónde está y dónde no está la
investigación. \textcite{zawacki2019systematic} mapearon 146 artículos sobre
aplicaciones de IA en educación superior: solo 13 (8,9\,\%) tienen como primer autor
a alguien de Educación, el 73,3\,\% emplea métodos cuantitativos --- sobre todo
cuasiexperimentales, con estudios piloto que no controlan el efecto novedad ---, y
solo 2 de 146 (1,4\,\%) reflexionan críticamente sobre implicaciones éticas. En su
tipología, la curación de materiales aparece en 3 estudios y el apoyo al docente en
el diseño de la enseñanza, en otros 3.

\textcite{bond2024meta} confirman el patrón en una revisión terciaria de 66
revisiones centradas en educación superior (cribado doble con $\kappa = 0{,}89$ y
$\kappa = 0{,}85$): solo 5 revisiones (7,6\,\%) mencionan el uso de IA para evaluar
materiales de aprendizaje --- y casi siempre midiendo tiempo de uso en plataforma,
no el contenido ---, y la curación de materiales aparece en 3 (4,5\,\%). En cuanto a
rigor, la puntuación media de calidad metodológica fue de 6,57 sobre 10 y
\textcquote[p.~34]{bond2024meta}{A noticeable 65\% of reviews are
critically low to medium quality}; el 51,5\,\% no reporta acuerdo entre revisores.
Entre los desafíos más citados figura el \textbf{desplazamiento de la autoridad}
hacia la IA (22,6\,\%), y entre los vacíos, la necesidad de mayor rigor y variedad
metodológica (36,4\,\%) junto con el llamado a evaluar el efecto de las herramientas
sobre el aprendizaje y las personas, y no solo su exactitud técnica.

\begin{advertencia}
Ninguna de estas dos revisiones evalúa la eficacia de herramienta alguna: son mapeos
de literatura. Que la curación de materiales aparezca en pocas revisiones indica
escasa \emph{síntesis} sobre el tema, no necesariamente escasa investigación
primaria. La brecha que se formula en la Parte~II se apoya en esa evidencia con esa
salvedad explícita.
\end{advertencia}
